\documentclass[11pt]{article}

\usepackage[a4paper,margin=1in]{geometry}
\usepackage{amsmath,amssymb,amsthm,mathtools}
\usepackage{graphicx}
\usepackage{microtype}
\usepackage{xcolor}
\usepackage{enumitem}
\usepackage[hidelinks]{hyperref}

\definecolor{statusblue}{RGB}{32,78,121}
\newtheorem{theorem}{Theorem}
\newtheorem{lemma}{Lemma}
\newcommand{\R}{\mathbb R}
\newcommand{\B}{\mathcal B}
\newcommand{\norm}[1]{\left\lVert#1\right\rVert}
\newcommand{\Sp}{\operatorname{Sp}}
\newcommand{\Spu}{\operatorname{Sp}_{u}}

\title{\textbf{A Spectral-Synthesis-Free Katznelson--Tzafriri Theorem for Measures}\\
\large A candidate full answer to Remark 1.3 of arXiv:1410.1301}
\author{Automated mathematics run \texttt{fa\_banach\_001}\\
Lane 14, model GPT5.6}
\date{9 August 2026}

\begin{document}
\maketitle

\begin{abstract}
Seifert asks whether the spectral-synthesis assumption can be removed from
the Hilbert-space Katznelson--Tzafriri theorem for finite measures. We give
an affirmative answer. The high-frequency resolvent hypothesis forces
asymptotic uniform norm continuity. This supplies exactly the regularity
needed to combine an \(L^1\)-smoothed Hilbert ultrapower with the unitary
asymptote and the spectral theorem, even when the original measure is
singular.
\end{abstract}

\noindent
\colorbox{statusblue!10}{\parbox{0.96\linewidth}{
\textbf{Status.} Candidate full result; verifier verdict: likely valid.
No spectral-synthesis assumption remains. Human review is recommended,
especially for the continuous-ultrapower and measure-integral steps.
}}

\section{Source question}

The source is Seifert, \emph{A Katznelson--Tzafriri theorem for measures},
arXiv:1410.1301 \cite{SeifertMeasures}. For a finite complex Borel measure
\(\mu\) on \(\R_+\), it uses
\[
 (\mathcal F\mu)(s)=\int_{\R_+}e^{-ist}\,d\mu(t),
 \qquad
 \widehat\mu(T)x=\int_{\R_+}T(t)x\,d\mu(t).
\]
Theorem 1.2 assumes that \(T\) is a bounded \(C_0\)-semigroup on a complex
Hilbert space, with generator \(A\), and that for some \(R>0\),
\begin{equation}\label{eq:resolvent-hyp}
 \{is:|s|\ge R\}\subset\rho(A),
 \qquad
 \sup_{|s|\ge R}\norm{R(is,A)}<\infty.
\end{equation}
If \(i\sigma(A)\cap\R\) is additionally a set of spectral synthesis and
\(\mathcal F\mu\) vanishes there, it concludes that
\(\norm{T(t)\widehat\mu(T)}\to0\).

Remark 1.3 on p.~2 asks:
\begin{quote}
\emph{It remains open whether [the spectral-synthesis assumption can be
dropped] also in the setting of Theorem 1.2.}
\end{quote}

\begin{figure}[ht]
\centering
\includegraphics[width=\textwidth]{figures/open_problem_crop.png}
\caption{Theorem 1.2 and Remark 1.3 on p.~2 of arXiv:1410.1301.}
\end{figure}

\section{Proof intuition}

For \(L^1\) kernels, Seifert's Hilbert-space proof smooths the functional
calculus in variation norm before passing to an ultrapower
\cite{SeifertHilbert}. That cannot be copied literally for a singular
measure: convolution with an \(L^1\) approximate identity does not converge
to a singular measure in total variation.

The resolvent assumption provides the missing replacement. On Hilbert space
it gives a negative critical growth bound, which implies that \(T\) becomes
uniformly norm-continuous far out on the time axis. We therefore smooth the
\emph{almost-norming vectors} rather than the measure. The smoothing places
those vectors in the continuous part of a Hilbert ultrapower, while tail
norm-continuity preserves the non-decay witnessed by the original singular
measure.

The unitary asymptote of that continuous ultrapower sees only the original
boundary spectrum. On this unitary group the spectral theorem turns the
measure calculus into multiplication by \(\mathcal F\mu\), so the induced
operator is zero. The resulting orbit must both decay and stay uniformly
away from zero, which is the contradiction.

\section{Main result}

\begin{theorem}\label{thm:main}
Let \(H\) be a complex Hilbert space and let \(T=(T(t))_{t\ge0}\) be a
bounded \(C_0\)-semigroup on \(H\), with generator \(A\), satisfying
\eqref{eq:resolvent-hyp}. Let \(\mu\) be a finite complex Borel measure on
\(\R_+\) such that
\[
 (\mathcal F\mu)(s)=0
 \qquad\bigl(s\in i\sigma(A)\cap\R\bigr).
\]
Then
\[
 \lim_{t\to\infty}\norm{T(t)\widehat\mu(T)}=0.
\]
In particular, spectral synthesis is unnecessary in Theorem 1.2 of
\cite{SeifertMeasures}.
\end{theorem}

\section{Tail norm continuity}

Recall the critical growth bound \(\delta(T)\): it is the infimum of the
growth bounds of \(T-V\) over all norm-continuous
\(V:(0,\infty)\to\B(H)\). The non-analytic Gearhart--Pr\"uss theorem and the
standard inequalities for these growth bounds give, on Hilbert space,
\[
 \delta(T)=s_0^\infty(A).
\]
This equality and its use under precisely \eqref{eq:resolvent-hyp} are
recorded in Sections 3--4 of \cite{SeifertMeasures}, drawing on
\cite{BBS,BS}. A Neumann-series argument applied to the uniformly bounded
resolvents in \eqref{eq:resolvent-hyp} gives
\begin{equation}\label{eq:delta-neg}
 \delta(T)=s_0^\infty(A)<0.
\end{equation}

\begin{lemma}[asymptotic uniform norm continuity]\label{lem:tail-cont}
If a bounded \(C_0\)-semigroup satisfies \(\delta(T)<0\), then for every
\(\gamma>0\) there are \(r_0>0\) and \(\eta>0\) such that
\[
 \norm{T(r+h)-T(r)}<\gamma
\]
whenever \(r\ge r_0\), \(|h|\le\eta\), and \(r+h\ge0\).
\end{lemma}

\begin{proof}
Choose a norm-continuous \(V:(0,\infty)\to\B(H)\), \(C>0\), and \(c>0\)
such that
\[
 \norm{T(r)-V(r)}\le Ce^{-cr}\qquad(r>0).
\]
For
\[
 D(r)=\limsup_{h\downarrow0}\norm{T(r+h)-T(r)},
\]
norm continuity of \(V\) gives \(D(r)\le2Ce^{-cr}\), hence \(D(r)\to0\).

Put \(M=\sup_{t\ge0}\norm{T(t)}\). Choose \(r_1\) so large that
\(MD(r_1)<\gamma/2\). By the definition of the limsup, after decreasing
\(\eta>0\) if necessary,
\[
 \norm{T(r_1+h)-T(r_1)}<\gamma/M
 \qquad(0\le h\le\eta).
\]
For \(r=r_1+q\), \(q\ge0\), the semigroup law gives
\[
 \norm{T(r+h)-T(r)}
 \le M\norm{T(r_1+h)-T(r_1)}<\gamma.
\]
This handles \(h\ge0\). If \(-\eta\le h<0\) and \(r\ge r_1+\eta\), apply
the positive-increment estimate at the base point \(r+h\). Thus
\(r_0=r_1+\eta\) works.
\end{proof}

\section{The continuous ultrapower}

We isolate the standard ultrapower facts used in the proof. Let
\(\mathcal U\) be a free ultrafilter on \(\mathbb N\), and let
\(H_{\mathcal U}\) be the Hilbert ultrapower. The operators
\[
 \widetilde T(t)[(x_n)]=[(T(t)x_n)]
\]
form a bounded semigroup, but need not be strongly continuous on all of
\(H_{\mathcal U}\). For \(a\in L^1(\R_+)\), let
\(\widehat a(\widetilde T)\) mean the coordinatewise ultrapower of
\(\widehat a(T)\), and define
\begin{equation}\label{eq:Hc}
 H_c=\overline{\operatorname{span}}
 \{\widehat a(\widetilde T)\xi:a\in L^1(\R_+),\ \xi\in H_{\mathcal U}\}.
\end{equation}

\begin{lemma}\label{lem:ultrapower}
The restriction \(S(t)=\widetilde T(t)|_{H_c}\) is a bounded \(C_0\)-semigroup.
For every \(a\in L^1(\R_+)\),
\begin{equation}\label{eq:integrated-bound}
 \widehat a(S)=\widehat a(\widetilde T)|_{H_c},
 \qquad
 \norm{\widehat a(S)}\le\norm{\widehat a(T)}.
\end{equation}
Consequently
\begin{equation}\label{eq:spectrum-inclusion}
 \Spu(S)\subset\Spu(T).
\end{equation}
\end{lemma}

\begin{proof}
If \(a\in L^1(\R_+)\), translation continuity in \(L^1\) gives
\[
 \norm{T(h)\widehat a(T)-\widehat a(T)}\longrightarrow0
 \qquad(h\downarrow0).
\]
The same estimate holds in the ultrapower norm. Hence \(\widetilde T\) is
strongly continuous on the generators in \eqref{eq:Hc}, and therefore on
their closed span. Invariance is immediate.

On a generator \(\widehat b(\widetilde T)\xi\), Fubini's theorem gives
\[
 \widehat a(S)\widehat b(\widetilde T)\xi
 =\widehat{a*b}(\widetilde T)\xi
 =\widehat a(\widetilde T)\widehat b(\widetilde T)\xi.
\]
This proves the first identity in \eqref{eq:integrated-bound} by density.
The coordinatewise definition gives the norm inequality.

For completeness, the spectrum of a bounded semigroup representation is
characterized by
\[
 \Sp(T)=\{\chi:|\widehat a(\chi)|\le\norm{\widehat a(T)}
 \text{ for every }a\in L^1(\R_+)\}.
\]
Thus \eqref{eq:integrated-bound} implies \(\Sp(S)\subset\Sp(T)\), and
intersecting with the unitary characters proves
\eqref{eq:spectrum-inclusion}. This is the same spectrum formalism used in
\cite{SeifertHilbert}.
\end{proof}

We shall also use the standard Hilbert-space unitary-asymptote construction
from the proof of Proposition 3.4 of \cite{SeifertHilbert}.

\begin{lemma}[unitary asymptote]\label{lem:asymptote}
For a bounded \(C_0\)-semigroup \(S\) on a Hilbert space \(G\), there exist a
Hilbert space \(K\), a strongly continuous unitary group \(U=(U(t))_{t\in\R}\)
on \(K\), and a bounded operator \(\pi:G\to K\) such that
\begin{align}
 \pi S(t)&=U(t)\pi \qquad(t\ge0),\label{eq:intertwine}\\
 \ker\pi&=\{x\in G:\norm{S(t)x}\to0\},\label{eq:kernel}\\
 \Sp(U)&\subset\Spu(S).\label{eq:unitary-spectrum}
\end{align}
\end{lemma}

\begin{proof}[Construction reminder]
Apply an invariant Banach limit to
\(t\mapsto\langle S(t)x,S(t)y\rangle\), quotient by the resulting null
space, and complete. This produces an isometric semigroup and the
intertwiner. Boundedness of \(S\) shows that the null space is exactly the
stable vectors: once an orbit has a subnet tending to zero, all of its tail
does. The minimal unitary extension gives \(U\). The integrated-form
spectrum characterization gives \eqref{eq:unitary-spectrum}; full details
are supplied in the cited unitary-asymptote construction
\cite[proof of Proposition 3.4]{SeifertHilbert}.
\end{proof}

\section{Proof of Theorem \ref{thm:main}}

Put
\[
 M=\sup_{t\ge0}\norm{T(t)},
 \qquad B=\widehat\mu(T).
\]
Suppose for contradiction that \(\norm{T(t)B}\not\to0\). Then there are
\(\varepsilon>0\), \(t_n\to\infty\), and unit vectors \(y_n\in H\) such
that
\begin{equation}\label{eq:failure}
 \norm{T(t_n)By_n}\ge\varepsilon.
\end{equation}
For \(0\le q\le t_n\), the semigroup law and boundedness imply
\begin{equation}\label{eq:backward-lower}
 \norm{T(q)By_n}\ge\varepsilon/M,
\end{equation}
because \(T(t_n)B=T(t_n-q)T(q)B\).

Fix \(\tau>0\). Since failure implies \(B\ne0\), Lemma
\ref{lem:tail-cont} supplies \(\eta>0\) and \(s_0>0\) such that
\begin{equation}\label{eq:tail-choice}
 \sup_{0\le h\le\eta}
 \norm{T(s+\tau+h)-T(s+\tau)}
 <\frac{\varepsilon}{2M\norm B}
 \qquad(s\ge s_0).
\end{equation}
Choose a nonnegative \(b\in L^1(\R_+)\), supported in
\([\tau,\tau+\eta]\), with \(\int b=1\).

Use Lemma \ref{lem:ultrapower}, put \(y=[(y_n)]\in H_{\mathcal U}\), and set
\[
 x=\widehat b(\widetilde T)y
   =[(\widehat b(T)y_n)]\in H_c.
\]
The map \(q\mapsto T(q)\widehat b(T)\) is norm-continuous. Indeed, for
positive \(h\), translation continuity of \(b\) gives
\[
 \sup_{q\ge0}\norm{T(q+h)\widehat b(T)-T(q)\widehat b(T)}
 \le M\norm{T(h)\widehat b(T)-\widehat b(T)}\longrightarrow0.
\]
It follows that this bounded norm-continuous operator-valued map is Bochner
integrable against \(|\mu|\). Approximation in
\(L^1(|\mu|;\B(H))\) by common simple functions gives operator-norm
convergence before passing to the ultrapower. Hence, using Fubini,
\begin{equation}\label{eq:measure-passes}
 z:=\widehat\mu(S)x
 =[(B\widehat b(T)y_n)]
 =[\widehat{\mu*b}(T)y_n]\in H_c.
\end{equation}
This is the point at which the \(L^1\) smoothing of \(x\) is essential; no
variation-norm approximation of \(\mu\) is asserted.

For \(s\ge s_0\), equations \eqref{eq:backward-lower} and
\eqref{eq:tail-choice} give
\begin{align}
 \norm{S(s)z}
 &=\lim_{\mathcal U}
   \norm{T(s)B\widehat b(T)y_n}\notag\\
 &=\lim_{\mathcal U}
   \norm{\int_{\tau}^{\tau+\eta}b(q)T(s+q)By_n\,dq}\notag\\
 &\ge\lim_{\mathcal U}\norm{T(s+\tau)By_n}
      -\frac{\varepsilon}{2M}
 \ge\frac{\varepsilon}{2M}.\label{eq:nondecay}
\end{align}
For the last inequality, the cofinite set on which \(t_n\ge s+\tau\)
belongs to \(\mathcal U\).

Apply Lemma \ref{lem:asymptote} to \(S\). From
\eqref{eq:intertwine}, bounded linearity of \(\pi\), and Bochner
integration,
\begin{equation}\label{eq:measure-intertwine}
 \pi z=\int_{\R_+}U(q)\pi x\,d\mu(q).
\end{equation}
Under the standard identification of unitary characters with real Fourier
parameters, \eqref{eq:spectrum-inclusion} and
\eqref{eq:unitary-spectrum} place the spectral measure of \(U\) inside
\(i\sigma(A)\cap\R\). (Explicitly, the character \(t\mapsto e^{-ist}\)
has generator \(-is\), so \(s=i(-is)\).) The spectral theorem and the
hypothesis on \(\mathcal F\mu\) therefore give
\[
 \int_{\R_+}U(q)\,d\mu(q)=0.
\]
Thus \(\pi z=0\). By \eqref{eq:kernel},
\(\norm{S(s)z}\to0\), contradicting \eqref{eq:nondecay}. The theorem is
proved.

\section{Verification, scope, and novelty}

\begin{itemize}[leftmargin=1.5em]
\item The proof never approximates a singular measure in total variation.
  The only smoothing is \(x=\widehat b(\widetilde T)y\), which makes the
  operator orbit norm-continuous and justifies
  \eqref{eq:measure-passes}.
\item The ultrapower is restricted to the \(L^1\)-essential subspace
  \eqref{eq:Hc}; strong continuity is not claimed on the full ultrapower.
\item No sequential compactness of an operator ball is used. The free
  ultrafilter preserves the almost-norming sequence directly.
\item The resolvent hypothesis is used only to obtain
  \(\delta(T)<0\), hence Lemma \ref{lem:tail-cont}. The argument proves the
  slightly more general statement with \(\delta(T)<0\) in place of
  \eqref{eq:resolvent-hyp}.
\item The result removes spectral synthesis only in the Hilbert-space
  theorem asked about in Remark 1.3. It makes no analogous assertion for
  arbitrary Banach spaces.
\end{itemize}

\paragraph{Novelty check.}
On 9 August 2026, bounded web/arXiv searches used the exact sentence in
Remark 1.3, arXiv:1410.1301, the title and author, and combinations of
\emph{bounded Borel measure}, \emph{Hilbert}, \emph{spectral synthesis}, and
\emph{Katznelson--Tzafriri}. They found the source, the companion Hilbert
paper \cite{SeifertHilbert}, the original measure theorem
\cite{BCT}, and the 2022 survey \cite{BattySeifertSurvey}, but no later work
claiming this removal. This is a bounded novelty check, not a comprehensive
priority search.

\paragraph{Human review recommendation.}
Check Lemma \ref{lem:tail-cont}, the integrated-form spectrum inclusion in
Lemma \ref{lem:ultrapower}, the uniform operator-norm approximation behind
\eqref{eq:measure-passes}, and the Fourier-sign convention in the final
spectral-theorem step. The verifier verdict is \emph{likely valid}.

\begin{thebibliography}{9}

\bibitem{SeifertMeasures}
D. Seifert,
\emph{A Katznelson--Tzafriri theorem for measures},
Integral Equations Operator Theory 81 (2015), 255--270;
arXiv:1410.1301.

\bibitem{SeifertHilbert}
D. Seifert,
\emph{Some improvements of the Katznelson--Tzafriri theorem on Hilbert
space}, Proc. Amer. Math. Soc. 143 (2015), 3827--3838;
arXiv:1410.1294.

\bibitem{BCT}
C. J. K. Batty, R. Chill, and Y. Tomilov,
\emph{Fine scales of decay of operator semigroups},
J. Eur. Math. Soc. 18 (2016), 853--929.

\bibitem{BBS}
C. J. K. Batty, M. D. Blake, and S. Srivastava,
\emph{A non-analytic growth bound for Laplace transforms and semigroups of
operators}, Integral Equations Operator Theory 45 (2003), 125--154.

\bibitem{BS}
C. J. K. Batty and S. Srivastava,
\emph{The non-analytic growth bound of a \(C_0\)-semigroup and inhomogeneous
Cauchy problems}, J. Differential Equations 194 (2003), 300--327.

\bibitem{BattySeifertSurvey}
C. J. K. Batty and D. Seifert,
\emph{Some developments around the Katznelson--Tzafriri theorem},
Acta Sci. Math. (Szeged) 88 (2022), 53--84; arXiv:2204.13411.

\end{thebibliography}

\end{document}
